\documentclass[12pt]{article}
\usepackage[a4paper,margin=1in]{geometry}
\usepackage[T1]{fontenc}
\usepackage[utf8]{inputenc}
\usepackage{setspace}
\setstretch{1.15}

\title{MaternaGuard Final Demo Script (2 Minutes)}
\author{Team Kundanzz}
\date{}

\begin{document}
\maketitle

Hello everyone, we are Team Kundanzz.
Before we begin, MaternaGuard is now live and deployed: frontend on Vercel and backend on Render.

I am Anirudh, working on the app and user workflow.
Aditya leads the ML pipeline and model optimization.
Donil handles backend APIs, database integration, and deployment.

MaternaGuard is built to reduce preventable maternal risk escalation delays in rural India.
The user enters six basic vitals available at PHC level: age, systolic BP, diastolic BP, blood sugar, body temperature, and heart rate.

Now I will show the demo flow.
On the Log Health page, we enter these six values for a patient.
When we submit, the app calls our FastAPI predict endpoint.
The backend runs our integrated ensemble inference pipeline and returns one of three classes: Low Risk, Mid Risk, or High Risk.
It also returns confidence and the top contributing factor for explainability.

On the dashboard, we can see recent assessments and risk trends.
In alerts, high-risk entries are clearly highlighted so frontline workers can prioritize escalation.
On the provider side, doctors can view patient-level assessment history to support faster triage decisions.

From the engineering side, we implemented modular backend routes for health, predict, sync, patients, and auth, with persistent storage and deployment-ready configuration.
Our ML stack combines multiple trained models and feature fusion to improve robustness on real-world variability.

MaternaGuard is designed for low-resource environments, with practical workflows for constrained connectivity.
Our next milestones are broader validation and stronger production deployment.

Thank you.

\end{document}
